\section{Proofs for the Online Residual and Tracking Results}\label{app:online}

\subsection{Normalized projected ascent}

\begin{lemma}[One-step normalized projected ascent]\label{lem:projected-ascent}
Let $J$ be $L$-smooth on the nonempty closed convex set $\Theta$.  Fix $x\in\Theta$, let $g=\nabla J(x)$ with $\norm{g}_2\le B_g$, and let $\widehat g$ be any finite vector.  Set
\[
 x^+=\Pi_\Theta\{x+\gamma\mathsf d(\widehat g)\},
 \qquad
 Q(x)=\gamma^{-1}[\Pi_\Theta\{x+\gamma\mathsf d(g)\}-x].
\]
If $\gamma L\le\vartheta$ and $C_g:=\max\{B_g,\vartheta\}$, then
\begin{equation}\label{eq:one-step-ascent}
 J(x^+)-J(x)
 \ge\frac{\vartheta\gamma}{8}\norm{Q(x)}_2^2
 -\frac{5C_g\gamma}{4\vartheta^2}\norm{\widehat g-g}_2^2.
\end{equation}
\end{lemma}

\begin{proof}[Proof of \Cref{lem:projected-ascent}]
The target is a lower bound in terms of the population projected residual $Q(x)$.  We first obtain a bound in terms of the actual displacement produced by $\widehat g$, then compare that displacement with the population displacement.

Let
\[
 \Delta x=x^+-x,
 \qquad
 c_g=\max\{\norm{g}_2,\vartheta\},
 \qquad
 \mathfrak e=\mathsf d(\widehat g)-\mathsf d(g).
\]
The projection variational inequality with the feasible comparison point $x$ gives
\begin{equation}\label{eq:projection-vi}
 \langle\mathsf d(\widehat g),\Delta x\rangle\ge\frac{\norm{\Delta x}_2^2}{\gamma}.
\end{equation}
Since $g=c_g\mathsf d(g)$, the $L$-smooth lower quadratic bound and \eqref{eq:projection-vi} imply
\begin{align}
 J(x+\Delta x)-J(x)
 &\ge \langle g,\Delta x\rangle-\frac L2\norm{\Delta x}_2^2\notag\\
 &=c_g\langle\mathsf d(\widehat g)-\mathfrak e,\Delta x\rangle
 -\frac L2\norm{\Delta x}_2^2\notag\\
 &\ge\left(\frac{c_g}{\gamma}-\frac L2\right)\norm{\Delta x}_2^2
 -c_g\norm{\mathfrak e}_2\norm{\Delta x}_2\notag\\
 &\ge\left(\frac{3c_g}{4\gamma}-\frac L2\right)\norm{\Delta x}_2^2
 -c_g\gamma\norm{\mathfrak e}_2^2\notag\\
 &\ge\frac{c_g}{4\gamma}\norm{\Delta x}_2^2
 -c_g\gamma\norm{\mathfrak e}_2^2.
 \label{eq:actual-displacement-ascent}
\end{align}
The penultimate line uses
$uv\le u^2/(4\gamma)+\gamma v^2$, and the last line uses
$\gamma L\le \vartheta\le c_g$.

Now let
\[
 t=\Pi_\Theta\{x+\gamma\mathsf d(g)\}-x=\gamma Q(x).
\]
Nonexpansiveness of Euclidean projection gives
$\norm{\Delta x-t}_2\le\gamma\norm{\mathfrak e}_2$.  From
$\norm{t}_2^2\le2\norm{\Delta x}_2^2+2\norm{\Delta x-t}_2^2$,
\[
 \norm{\Delta x}_2^2\ge\frac12\norm{t}_2^2-\gamma^2\norm{\mathfrak e}_2^2.
\]
Substituting into \eqref{eq:actual-displacement-ascent} yields
\[
 J(x^+)-J(x)
 \ge\frac{c_g\gamma}{8}\norm{Q(x)}_2^2
 -\frac{5c_g\gamma}{4}\norm{\mathfrak e}_2^2.
\]
Finally,
$\mathsf d(g)=\Pi_{\{u:\norm{u}_2\le1\}}(g/\vartheta)$, so nonexpansiveness of the unit-ball projection implies
\[
 \norm{\mathfrak e}_2\le \vartheta^{-1}\norm{\widehat g-g}_2.
\]
Using $c_g\ge \vartheta$ in the positive term and
$c_g\le C_g=\max(B_g,\vartheta)$ in the negative term proves \eqref{eq:one-step-ascent}.
\end{proof}

\subsection{Average expected squared projected residual}

\begin{proof}[Proof of \Cref{thm:online-residual}]
We first connect the final estimator to the true gradient.  Conditional on $\cF_k$, \Cref{thm:estimator} gives
\[
 \E(\widehat g_k^{\db}-\widetilde g_k(\theta_k)\mid\cF_k)=0,
 \qquad
 \E(\norm{\widehat g_k^{\db}-\widetilde g_k(\theta_k)}_2^2\mid\cF_k)
 \le\frac{C_{\rm var}(n)}{M_k}.
\]
The cross term with the deterministic conditional vector
$\widetilde g_k(\theta_k)-g_k(\theta_k)$ therefore vanishes.  Hence
\begin{align*}
 \E\norm{\widehat g_k^{\db}-g_k(\theta_k)}_2^2
 &=\E\norm{\widehat g_k^{\db}-\widetilde g_k(\theta_k)}_2^2
 +\E\norm{\widetilde g_k(\theta_k)-g_k(\theta_k)}_2^2\\
 &\le\frac{C_{\rm var}(n)}{M_k}+\E\beta_k^2,
\end{align*}
which is \eqref{eq:true-gradient-mse}.

We next apply \Cref{lem:projected-ascent} to the true episode objective $J_k$ at $x=\theta_k$ and $\widehat g=\widehat g_k^{\db}$.  Let
\[
 \mathfrak j_k=J_k(\theta_{k+1})-J_k(\theta_k).
\]
After taking expectations,
\begin{equation}\label{eq:online-one-step}
 \frac{\vartheta\gamma}{8}\E\zeta_k
 \le\E \mathfrak j_k
 +\frac{5C_g\gamma}{4\vartheta^2}
 \E\norm{\widehat g_k^{\db}-g_k(\theta_k)}_2^2.
\end{equation}
The objective increments telescope across moving objectives as the pathwise identity
\begin{align}\label{eq:moving-objective-telescope}
 \sum_{k=1}^{K}\mathfrak j_k
 ={}&J_K(\theta_{K+1})-J_1(\theta_1)\notag\\
 &+\sum_{k=1}^{K-1}
 \{J_k(\theta_{k+1})-J_{k+1}(\theta_{k+1})\}.
\end{align}
Since $\abs{J_k}\le B_v$ and the second line is bounded above by $\cV_K$,
\begin{equation}\label{eq:telescope-upper}
 \sum_{k=1}^{K}\mathfrak j_k\le2B_v+\cV_K.
\end{equation}
Summing \eqref{eq:online-one-step}, applying \eqref{eq:true-gradient-mse} and \eqref{eq:telescope-upper}, and dividing by $\vartheta\gamma K/8$ gives exactly \eqref{eq:online-residual-bound}.
\end{proof}

\subsection{Discounted dynamic local regret}

\begin{proof}[Proof of \Cref{cor:dynamic-local-regret}]
Write $W=W_{\rho,w}$.  With the convention $Q_t(\theta_t)=0$ for $t\le0$, the coefficients $\rho^j/W$, $j=0,\ldots,w-1$, are nonnegative and sum to one.  Convexity of $q\mapsto\norm{q}_2^2$ therefore gives
\begin{align*}
 \norm{\bar Q_k}_2^2
 &\le \frac1W\sum_{j=0}^{w-1}\rho^j
       \norm{Q_{k-j}(\theta_{k-j})}_2^2.
\end{align*}
Summing over $k\le K$ and reindexing,
\begin{align*}
 \Reg_{\rho,w}(K)
 &\le
 \sum_{t=1}^{K}\zeta_t
 \left\{\frac1W\sum_{j=0}^{\min\{w-1,K-t\}}\rho^j\right\}
 \le \sum_{t=1}^{K}\zeta_t,
\end{align*}
which proves \eqref{eq:dynamic-local-regret-bound}.  Taking expectations and dividing by $K$ gives \eqref{eq:average-dynamic-local-regret-bound}.
\end{proof}

\subsection{Stationary-set consequences}

\begin{lemma}[Drift of stationary sets]\label{lem:stationary-set-drift}
If \eqref{eq:error-bound} holds, define
\begin{equation}\label{eq:gradient-drift}
 \delta_k^g:=\sup_{\theta\in\Theta}
 \norm{g_{k+1}(\theta)-g_k(\theta)}_2.
\end{equation}
Then
\begin{equation}\label{eq:stationary-set-drift}
 \dH(\Sta_{k+1},\Sta_k)
 \le\frac{\kappa}{\vartheta}\delta_k^g.
\end{equation}
\end{lemma}

\begin{proof}[Proof of \Cref{cor:stationary-tracking}]
Condition \eqref{eq:error-bound} gives pointwise
\[
 \dist^2(\theta_k,\Sta_k)
 \le\kappa^2\norm{Q_k(\theta_k)}_2^2
 =\kappa^2\zeta_k.
\]
Taking expectations and averaging over $k$ proves \eqref{eq:stationary-tracking}.  Substitution of \eqref{eq:online-residual-bound} yields the stated finite-horizon control.
\end{proof}

\begin{proof}[Proof of \Cref{lem:stationary-set-drift}]
The normalized direction map is $\vartheta^{-1}$-Lipschitz because it is the Euclidean projection of $g/\vartheta$ onto the closed unit ball.  Using nonexpansiveness of $\Pi_\Theta$ in \eqref{eq:projected-residual},
\begin{align*}
 \sup_{\theta\in\Theta}
 \norm{Q_{k+1}(\theta)-Q_k(\theta)}_2
 &\le
 \sup_{\theta\in\Theta}
 \norm{\mathsf d(g_{k+1}(\theta))-\mathsf d(g_k(\theta))}_2\\
 &\le\frac{\delta_k^g}{\vartheta}.
\end{align*}
If $\theta\in\Sta_k$, then $Q_k(\theta)=0$, so \eqref{eq:error-bound} applied to episode $k+1$ gives
\[
 \dist(\theta,\Sta_{k+1})
 \le\kappa\norm{Q_{k+1}(\theta)}_2
 \le\frac{\kappa\delta_k^g}{\vartheta}.
\]
Taking the supremum over $\theta\in\Sta_k$ gives one directed Hausdorff bound.  Interchanging $k$ and $k+1$ gives the reverse bound because the same uniform gradient-difference bound is symmetric.  Taking the maximum proves \eqref{eq:stationary-set-drift}.
\end{proof}

\begin{proof}[Proof of \Cref{cor:vanishing-tracking}]
In \eqref{eq:online-residual-bound}, the endpoint term is $O(K^{-1})$.  The assumption $\E\cV_K=o(K)$ makes the variation term vanish after division by $K$.  Since $C_{\rm var}(n)$ is fixed,
\[
 \frac1K\sum_{k=1}^{K}\left(\frac{C_{\rm var}(n)}{M_k}+\E\beta_k^2\right)\to0
\]
by the assumed Ces\`aro condition.  Therefore $\cE_K\to0$.  Equation \eqref{eq:average-dynamic-local-regret-bound} gives the DLR conclusion, and \eqref{eq:stationary-tracking} gives the stationary-set conclusion when the local condition \eqref{eq:error-bound} holds.
\end{proof}

\paragraph{Objective variation in the semi-synthetic environment.}
Write $z_k=(\log X_k,r_k,\omega_k)$ and $\boldsymbol\chi_k=(\mu_k,b_k,s_k)$ for the episode-start account state and return-law coefficients.  On the bounded analysis domain of Assumption~\ref{ass:trajectory}, the semi-synthetic return bound and transaction-cost rule keep each daily wealth multiplier uniformly positive and bounded.  The fixed horizon therefore bounds the raw change in log wealth; reference values and holdings are bounded by the same assumption and the portfolio simplex.  The finite factor-block library and compact parameter set keep Dirichlet shape parameters in a common compact subset of $(0,\infty)$.  Uniform likelihood moments, the Lipschitz executable-action map in preceding holdings, and the fixed-horizon wealth and reference recursions give uniform sensitivity of the conditional outcome law to $z$ and $\boldsymbol\chi$.  The regularized CPT value and probability-weight functions then yield constants $C_1,C_2<\infty$, independent of $k$ and $\theta$, such that
\[
 \sup_{\theta\in\Theta}\abs{J_{k+1}(\theta,r_{k+1})-J_k(\theta,r_k)}
 \le C_1\bigl(\abs{\log X_{k+1}-\log X_k}+\abs{r_{k+1}-r_k}
       +\norm{\omega_{k+1}-\omega_k}_1\bigr)
       +C_2\norm{\boldsymbol\chi_{k+1}-\boldsymbol\chi_k}.
\]
The episode-boundary transition gives $\norm{z_{k+1}-z_k}=O(\alpha_k)$.  The return-law coefficients change at one boundary, then along a finite interpolation, and are fixed thereafter; thus $\sum_k\norm{\boldsymbol\chi_{k+1}-\boldsymbol\chi_k}<\infty$.  Consequently $\cV_K=O(K^{1/4})=o(K)$ on the stated analysis domain.  The training simulator samples the same conditional factor-block and return law as the environment, so its population gradient error is zero under that model.
